\section{Experimental Protocol}
\label{sec:experiments}

\subsection{Common setup}

All experiments use the same frozen ACT checkpoint for LIBERO Object, with a
prediction horizon of 100 and seven-dimensional relative actions.  Dimensions
0--5 form the arm/end-effector group and dimension 6 the gripper group.  Control
runs at 20 Hz.  Temporal ensembling is disabled, no policy weights are updated,
and compared configurations use the same checkpoint and paired official initial
states.  Episode success is the benchmark success condition.  Policy-query rate
is the number of full policy evaluations divided by environment steps.

Each episode is seeded by 1000 plus its official initial-state index and stops on
benchmark success or the environment's 280-step limit.  The policy output is
checked at runtime to have shape $100\!\times\!7$.  We log every environment
step, policy query, group cursor, source-generation identifier, and source age;
the paper's reported numbers come from committed aggregates over these logs.
As an implementation check, all five task-0 diagonal group configurations
reproduce their corresponding global success, episode length, and query counts.

We report paired success counts and, for selected comparisons, a two-sided exact
McNemar diagnostic.  These small-sample diagnostics are descriptive; task-level
tests are not corrected for the many explored configurations.

\subsection{Full task-0 landscape}

For `pick up the alphabet soup and place it in the basket,' we evaluate official
initial states 0--49.  Global controls use
$h\in\{1,2,4,8,16\}$.  Group-specific execution evaluates the complete
$5\!\times\!5$ Cartesian product over the same horizon set, producing exact
diagonal controls and ten symmetric assignment pairs $(h_a,h_g)$ versus
$(h_g,h_a)$.  This protocol comprises 1,500 paired episodes.

\subsection{Cross-task diagnostic}

We next evaluate LIBERO Object tasks 1--9 on official states 0--19 per task,
with seeds 1000--1019.  The common coarse set contains diagonal/global-equivalent
horizons 2, 4, 8, and 16 and eight off-diagonal pairs:
$(2,8)$, $(2,16)$, $(8,2)$, $(8,16)$, $(16,2)$, $(16,8)$,
$(4,16)$, and $(16,4)$.  No task-specific policy training or tuning is performed.
This grid is deliberately coarser than the task-0 landscape, so comparisons of
per-task best points are retrospective over the evaluated set, not estimates of
a continuously optimized method.

Table~\ref{tab:protocols} summarizes the two protocols.  The 2,180 cross-task
episodes include a 20-episode global/diagonal duplicate used only to verify
equivalence.  Scientific comparisons use 12 unique static points per task.

\begin{table}[t]
\caption{Evaluation protocols}
\label{tab:protocols}
\centering
\scriptsize
\setlength{\tabcolsep}{3.0pt}
\begin{tabular}{lcccc}
\toprule
Protocol & Tasks & States/task & Unique points & Episodes \\
\midrule
Full landscape & 0 & 50 & 30 & 1,500 \\
Cross-task & 1--9 & 20 & 12 & 2,180 \\
\bottomrule
\end{tabular}
\end{table}
